\chapter{Conclusion}
\label{chap:conclusion}

\section{Conclusions Against the Hypotheses}
% Answer each of H1-H4 explicitly (rubric: conclusions relate to the
% original research questions/hypotheses).

This thesis asked whether LLM essay evaluators score valid
non-Western rhetoric lower, and whether a culturally plural panel
can narrow and flag that disparity. Table~\ref{tab:verdicts} answers
each hypothesis against the evidence of Chapter~4.

\begin{table}[htbp]
\centering
\caption[Verdicts on hypotheses H1--H4]{Verdicts on the four
hypotheses. All statistics from Chapter~4; $n$=100 per cohort
(E1/E2), $n$=70 held-out (E3); 95\% writer-cluster bootstrap CIs.
The cohort gap is an unadjusted disparity (cohorts not matched on
independently judged quality).}
\label{tab:verdicts}
\small
\begin{tabular}{p{2.2cm} p{2.2cm} p{7.6cm}}
\hline
Hypothesis & Verdict & Decisive evidence \\
\hline
H1 existence \& generality & Supported as an unadjusted
disparity &
Native$-$SEA gap positive on all four families, every CI excluding
zero: 1.18 (Claude), 0.80 (GPT-4o), 0.23 (Qwen3), 0.19 (Gemini)
bands ($d$ = 1.01, 1.09, 0.39, 0.43); larger for B1\_2 essays.
Cultural mechanism not isolated by the design \\[3pt]
H2 mitigation, $k{=}5 > k{=}3$ & Not supported (GPT-4o); rejected
(Claude) &
Direct test, adjudicated $k$=5 vs $k$=3: GPT-4o gap reduction 0.05
[$-0.03$, 0.14], association change $+0.02$ [$-0.06$, 0.12] (no
clear evidence); Claude gap reduction $-0.10$ [$-0.18$, $-0.02$]
and association change $-0.06$ [$-0.10$, $-0.02$] (worse on
both). Components: $k$=5 alone
added 0.02 [$-0.02$, 0.06] on GPT-4o and $-0.10$ [$-0.15$, $-0.06$]
on Claude; the $k$=3 panel had cut the gap 23--28\%, no more than
its SEA lens alone. Scope note: the $k$=3 panel amplified
the gap on Gemini (0.19 $\to$ 0.45), not carried into E2 \\[3pt]
H3 detection $\ge$90\%/$\le$5\% & Not supported &
At the specified FPR $\le$ 5\%: held-out sensitivity 0.056 (GPT-4o)
and 0.000 (Claude). Held-out AUC 0.639 and 0.481; $r$(BDS, human
dispersion) $=$ 0.23 and $-0.13$ \\[3pt]
H4 trade-off boundary & Not testable against graded truth; small
in-sample trade-off on the proxy (exploratory) &
Comparable graded ground truth is unavailable for the full cohort
sample (13 essays overlap the GRA), so accuracy is proxied by band
association. 10{,}626-vector sweep: equal weighting
dominated on both models (gap 0.55 $\to$ 0.32 with association 0.39
$\to$ 0.42 on GPT-4o; 1.01 $\to$ 0.79 with 0.35 $\to$ 0.36 on
Claude); along the non-dominated set higher association needs a
larger gap (0.32--0.38 / 0.42--0.43; 0.79--0.90 / 0.355--0.361);
minimum-gap vector is the SEA lens alone, by construction \\
\hline
\end{tabular}
\end{table}

The table is read as follows. An unadjusted native$-$SEA scoring
disparity exists on every model family tested, six-fold heterogeneous
in magnitude and larger for the intermediate-band essays; because the
cohorts were not matched on independently assessed quality, this
disparity is consistent with evaluative epistemic filtering but does
not establish it (H1). Culturally framed scoring prompts changed
the disparity in a model-dependent way, reducing it on the two
models where it was largest, by marking both cohorts down (the
native cohort more) and by no more than the Southeast Asian lens
alone achieved, and amplifying it on the model where it was
smallest; five lenses gave no clear further reduction on GPT-4o and
a larger gap on Claude, and the adjudicator no clear change (H2).
H4 could not be tested against graded ground truth; on the
proficiency-band proxy, equal weighting was inefficient and a small
in-sample trade-off remained along the non-dominated set (H4). The panel's
divergence signal did not meet its specified operating point, and
was larger on native than on Southeast Asian essays, so the
review-trigger use of BDS remains a hypothesis rather than a
validated capability (H3).

Turning to the research questions, RQ1 is answered in part: a
scoring disparity against Southeast Asian writing is present and
general, its size is model-specific, and isolating its cultural
cause requires matched human assessments this design did not
include. RQ2 is answered conditionally. Culturally
framed prompts can narrow the disparity at inference time on the
models where it is material, and panel scores are strongly associated with 80-rater human
consensus ($r \approx 0.9$; calibration was not assessed);
but the narrowing has not been shown to be an improvement in
fairness rather than a change of scoring criteria, and the
self-diagnostic signal is not yet reliable enough to route essays
autonomously.

\section{Contributions}
\label{sec:conclusion-contributions}

This thesis makes five contributions. Conceptually, it formalises
\emph{evaluative epistemic filtering} as a candidate bias mechanism
in LLM-driven scoring (Definition~\ref{def:eef}), grounded in
epistemic injustice and contrastive rhetoric. Architecturally, it specifies
CAMPS, a culturally-calibrated scoring panel with a divergence-based
human-review trigger, extending multi-perspective bias mitigation
from generation \citep{dorleon2026uncovering} to evaluation. Empirically, it evaluates the framework on topic-controlled ICNALE
cohorts across four model families, establishing the disparity at
scale and characterising where the panel changes it; and it
disconfirms two working assumptions, that adding the East Asian and
Mekong lenses would add mitigation and that inter-judge divergence
transfers directly into a detection signal. Methodologically, it introduces ROC calibration of
the Bias Divergence Score against the disagreement of eighty human
raters, an evaluation design that exposed the construct gap between
cultural contestation and general rater disagreement. Practically,
it gives deploying institutions a re-runnable baseline audit and an
account of what narrowing the disparity did and did not cost on
this evidence.

\section{Limitations}
\label{sec:limitations}

Beyond the validity threats analysed in Section~\ref{sec:threats},
four limitations bound this work. First, corpus scope: the evidence
rests on short, topic-controlled argumentative essays by university
students; high-stakes examination writing differs in genre, length,
and stakes. Second, the lenses cover five perspectives and leave others
(notably Middle Eastern, South Asian, and African traditions) to
future work; the constraint is the absence of validating cohorts and
documented rhetorical grounding for a lens, not of essays. Third,
CAMPS is an inference-time mitigation: it treats the symptom at the
point of judgement, not the training-distribution cause; a panel of
biased judges may be fairer than one, but is not unbiased. Fourth,
the framework multiplies inference cost by the panel size and is
intended to route contested essays to humans, a targeting E3 has
not validated; whether the cost is acceptable is an institutional
decision.

Stated plainly, the scope within which the findings hold is this.
CAMPS narrowed the disparity on the two frontier models whose
measured gaps were largest (0.8 bands or more, a description rather
than a threshold), on short argumentative essays by upper-band
Southeast Asian learners, with three lenses. It
had no clear effect on a model with a small gap and amplified the gap
on the model with the smallest. It should not be applied without first measuring the disparity it
is meant to treat. Five internal limitations qualify the numbers. The cohorts were not matched on
independently assessed quality, so the cohort gap is an unadjusted
disparity; the human-anchored check could compare only the learner
cohorts, since the rated archive holds four native essays. The lens
prompts add scoring criteria rather than reinterpreting a fixed
rubric, so the panel's effect may reflect changed expectations rather
than a corrected bias; a control panel of non-cultural prompt
variants, which would separate these explanations, was not run. BDS
was evaluated at band-level resolution (seven observed values). The
``accuracy'' axis of the weight sweep is association with a
two-level proficiency placement, not essay-level human quality. And
the 80-rater correlation of about 0.9 is association on a
single-topic set, not calibration or proof of a shared rubric.

\section{Future Work}
\label{sec:future-work}

The null and unresolved results define the nearest future work.
First, re-validating BDS with continuous, confidence-weighted lens
scores against annotations of specifically \emph{cultural}
contestation, so that the detector is tested on the construct it
measures, alongside the non-cultural control panel that would
isolate what cultural plurality adds over prompt averaging. Second,
graded human ratings on the cohort essays, so that
Definition~\ref{def:eef} can be tested as an LLM$-$human residual
between native and Southeast Asian writers rather than read from a
raw cohort gap. Beyond these: extending the panel to further rhetorical traditions
as validating corpora become available; a theoretical treatment, since
the formal guarantees for multi-perspective generation
\citep{dorleon2026uncovering} suggest an analogous proposition on
when panel consensus provably narrows cohort gaps; and a classroom
pilot of BDS-triggered review.
